\section{多重假设检验}
\begin{property}
	若控制每一次假设检验的显著性水平为$\alpha$，$n$次独立假设检验中至少犯一次第一次错误的概率为$1-(1-\alpha)^n$。
\end{property}
\begin{proof}
	此时一次第一类错误都没犯的概率为$(1-\alpha)^n$。
\end{proof}
设$n$个原假设中有$n_0$个为真，$n_0$个为假。记试验中被错误拒绝的假设个数为$V$，被正确拒绝的假设个数为$S$，总拒绝数为$R=V+S$。
\begin{definition}
	称$P(V\geqslant1)$为\gls{FWER}。
\end{definition}
\begin{definition}
	称：
	\begin{equation*}
		\operatorname{E}\left(\frac{V}{\max\{R,1\}}\right)
	\end{equation*}
	为\gls{FDR}。
\end{definition}
\begin{method}[BonFerroni]
	由\cref{prop:Measure}(3)（次有限可加性）可知，将$n$个假设检验问题中单个假设检验的显著性水平设置为$\dfrac{\alpha}{n}$，即可保证FWER小于$\alpha$。
\end{method}
\begin{method}[Benjamini--Hochberg]
	设$n$个假设检验对应的$p$值为$p_1,p_2\dots,p_n$，将其从小到大排序为$p_{(1)}\leqslant p_{(2)}\leqslant\cdots\leqslant p_{(n)}$。给定目标水平$\alpha$，取最大的$k$使得：
	\begin{equation*}
		p_{(k)}\leqslant\frac{k}{n}\alpha
	\end{equation*}
	则拒绝$H_{(1)},H_{(2)}\dots,H_{(k)}$。若各真原假设对应的$p$值独立且服从$U(0,1)$，并且与其余$p$值独立，则此时的FDR不超过$\alpha$。
\end{method}
\begin{proof}
	记$I_0$为真原假设的指标集合，$m_0=|I_0|$。对每个$i\in I_0$，记$R^{(-i)}$为删去$p_i$之后，其余$n-1$个$p$值按同一Benjamini--Hochberg规则所产生的拒绝数。由逐步上升法的结构可知，若第$i$个真原假设被错误拒绝且$R=r$，则必有
	\begin{equation*}
		p_i\leqslant \frac{r}{n}\alpha
	\end{equation*}
	并且
	\begin{equation*}
		R^{(-i)}=r-1
	\end{equation*}
	于是
	\begin{align*}
		\mathrm{FDR}
		&=E\left(\frac{V}{R\vee1}\right)\\
		&=\sum_{i\in I_0}E\left(\frac{\mathbf{1}\{p_i\text{被拒绝}\}}{R\vee1}\right)\\
		&=\sum_{i\in I_0}\sum_{r=1}^n\frac1rP\left(p_i\leqslant \frac{r}{n}\alpha,R=r\right)\\
		&\leqslant \sum_{i\in I_0}\sum_{r=1}^n\frac1rP\left(p_i\leqslant \frac{r}{n}\alpha,R^{(-i)}=r-1\right)
	\end{align*}
	若各真原假设对应的$p$值独立且服从$U(0,1)$，并且与其余$p$值独立，则
	\begin{align*}
		\mathrm{FDR}
		&\leqslant \sum_{i\in I_0}\sum_{r=1}^n\frac1rP\left(p_i\leqslant \frac{r}{n}\alpha\right)P\left(R^{(-i)}=r-1\right)\\
		&=\sum_{i\in I_0}\sum_{r=1}^n\frac1r\frac{r\alpha}{n}P\left(R^{(-i)}=r-1\right)\\
		&=\sum_{i\in I_0}\frac{\alpha}{n}\sum_{r=1}^nP\left(R^{(-i)}=r-1\right)\\
		&=\frac{m_0}{n}\alpha\\
		&\leqslant \alpha
	\end{align*}
	故Benjamini--Hochberg方法在独立情形下控制FDR
\end{proof}

\begin{method}[Benjamini--Yekutieli]
	设$n$个假设检验对应的$p$值为$p_1,\dots,p_n$，将其从小到大排序为
	\begin{equation*}
		p_{(1)}\leqslant p_{(2)}\leqslant\cdots\leqslant p_{(n)}
	\end{equation*}
	记
	\begin{equation*}
		c_n=\sum_{j=1}^n\frac1j
	\end{equation*}
	给定目标水平$\alpha$，取最大的$k$使得
	\begin{equation*}
		p_{(k)}\leqslant \frac{k}{n\,c_n}\alpha
	\end{equation*}
	则拒绝
	\begin{equation*}
		H_{(1)},\dots,H_{(k)}
	\end{equation*}
	证明：Benjamini--Yekutieli证明了如下上界：对任意依赖结构，若采用Benjamini--Hochberg型逐步上升法，并以
	\begin{equation*}
		\alpha_i=\frac{i}{n}\gamma,\qquad i=1,\dots,n
	\end{equation*}
	作为临界值，则有
	\begin{equation*}
		\mathrm{FDR}\leqslant \frac{m_0}{n}\gamma c_n
	\end{equation*}
	现在取
	\begin{equation*}
		\gamma=\frac{\alpha}{c_n}
	\end{equation*}
	则对应临界值恰为
	\begin{equation*}
		\alpha_i=\frac{i}{n\,c_n}\alpha
	\end{equation*}
	于是
	\begin{align*}
		\mathrm{FDR}
		&\leqslant \frac{m_0}{n}\frac{\alpha}{c_n}c_n\\
		&=\frac{m_0}{n}\alpha\\
		&\leqslant \alpha
	\end{align*}
	故Benjamini--Yekutieli方法在任意依赖结构下控制FDR
\end{method}